\chapter{The Symbolic Method}

The central ambition of analytic combinatorics is to pass from a structural description of a class of objects to an equation satisfied by its generating function, and then --- via the analytic machinery of Chapter~2 --- to asymptotic counting formulas. The \emph{symbolic method} is the first half of that program. It is a dictionary translating set-theoretic constructions on combinatorial classes into algebraic operations on generating functions.

There is one important warning at the outset. A symbolic specification is not just a vague resemblance between formulas and objects. It must describe each object \emph{exactly once}. If the same object can be built in two different ways, the generating function will overcount. So the symbolic method works best when a combinatorial decomposition is really a bijection between objects and structured tuples, sequences, or sets of simpler objects. The standard reference for the material in this chapter is \cite{FlajoletSedgewick2009}, Part~A.

\section{Combinatorial Classes}

\begin{definition}
A \emph{combinatorial class} is a pair $(\mathcal{C}, |\cdot|)$ where $\mathcal{C}$ is a set and
\[
  |\cdot| : \mathcal{C} \to \mathbb{Z}_{\ge 0}
\]
is a \emph{size function} such that, for every $n \ge 0$, the level set
\[
  \mathcal{C}_n = \{\gamma \in \mathcal{C} : |\gamma| = n\}
\]
is finite.
\end{definition}

Write $c_n = |\mathcal{C}_n|$. The \emph{ordinary generating function} (OGF) of $\mathcal{C}$ is
\[
  C(z) = \sum_{\gamma \in \mathcal{C}} z^{|\gamma|} = \sum_{n \ge 0} c_n z^n.
\]
The second equality comes from grouping together all objects of the same size $n$: there are exactly $c_n$ of them. The finiteness of each level set is what makes this a valid formal power series, because the coefficient of $z^n$ is the finite number $c_n$.

Three archetypal classes will recur throughout the book.
\begin{itemize}
  \item Binary words over $\{a,b\}$, with size equal to word length.
  \item Plane trees, where the children of each node are ordered from left to right.
  \item Permutations, which are naturally \emph{labeled} objects and therefore call for exponential generating functions rather than ordinary ones.
\end{itemize}
For example, the plane trees of sizes $1$ and $2$ look like a single root, and a root with one child; and a size-$3$ permutation is something like $(2,3,1)$, where the labels matter.

\begin{center}
\begin{tikzpicture}[
  level distance=0.6cm,
  sibling distance=0.8cm,
  every node/.style={circle, draw, fill=black!10, inner sep=1.5pt}
]
  \begin{scope}[shift={(0,0)}]
    \node[draw=none, fill=none, above=0.2cm] {Size 1};
    \node {};
  \end{scope}
  \begin{scope}[shift={(2.5,0)}]
    \node[draw=none, fill=none, above=0.2cm] {Size 2};
    \node {} child {node {}};
  \end{scope}
  \begin{scope}[shift={(5.5,0)}]
    \node[draw=none, fill=none, above=0.2cm] {Size 3};
    \node {} child {node {} child {node {}}};
  \end{scope}
  \begin{scope}[shift={(7.5,0)}]
    \node[draw=none, fill=none, above=0.2cm] {Size 3};
    \node {} child {node {}} child {node {}};
  \end{scope}
\end{tikzpicture}
\end{center}

Two tiny classes serve as universal building blocks.

\begin{definition}
The \emph{neutral class} $\mathcal{E} = \{\epsilon\}$ contains one object of size $0$, so its OGF is $E(z)=1$. The \emph{atomic class} $\mathcal{Z} = \{\bullet\}$ contains one object of size $1$, so its OGF is $Z(z)=z$.
\end{definition}

These are the combinatorial analogues of the constants $1$ and $z$. For instance, a word of length $3$ built from three atoms has the symbolic form
\[
  \mathcal{Z} \times \mathcal{Z} \times \mathcal{Z}.
\]
The rest of the symbolic method consists of combining $\mathcal{E}$ and $\mathcal{Z}$ through a small family of constructors.

\section{The Translation Dictionary for OGFs}

\begin{proposition}[Symbolic transfer, unlabeled case]
\label{prop:symbolic-ogf}
Let $\mathcal{A}$ and $\mathcal{B}$ be combinatorial classes with OGFs $A(z)$ and $B(z)$. Then the following constructions translate into the following OGF operations.
\begin{enumerate}
  \item \textbf{Disjoint union.} If every object is tagged as coming from the $\mathcal{A}$ side or the $\mathcal{B}$ side, then
  \[
    \mathcal{A} + \mathcal{B} \longleftrightarrow A(z)+B(z).
  \]
  \item \textbf{Cartesian product.} If objects are ordered pairs $(a,b)$ with size $|(a,b)|=|a|+|b|$, then
  \[
    \mathcal{A} \times \mathcal{B} \longleftrightarrow A(z)B(z).
  \]
  \item \textbf{Sequence.} If $\mathcal{A}$ contains no object of size $0$, then finite sequences of $\mathcal{A}$-objects satisfy
  \[
    \SEQ(\mathcal{A}) \longleftrightarrow \frac{1}{1-A(z)}.
  \]
\end{enumerate}
\end{proposition}

The word \emph{disjoint} matters. If the same underlying object could belong to both classes, it would be counted twice. When two classes are not literally disjoint, one uses a \emph{tagged disjoint union}: an object is paired with a label saying which side it came from.

\begin{proof}
For disjoint union, every object belongs to exactly one summand, so the OGF simply splits:
\[
  \sum_{\gamma \in \mathcal{A}+\mathcal{B}} z^{|\gamma|}
  = \sum_{\alpha \in \mathcal{A}} z^{|\alpha|} + \sum_{\beta \in \mathcal{B}} z^{|\beta|}
  = A(z)+B(z).
\]

For Cartesian products, we use the additive size rule:
\[
  \sum_{(a,b) \in \mathcal{A}\times\mathcal{B}} z^{|(a,b)|}
  = \sum_{(a,b)} z^{|a|+|b|}
  = \sum_{a \in \mathcal{A}} \sum_{b \in \mathcal{B}} z^{|a|}z^{|b|}
  = A(z)B(z).
\]
At the coefficient level,
\[
  [z^n]A(z)B(z) = \sum_{k=0}^n a_k b_{n-k},
\]
which has the expected combinatorial meaning: to build a size-$n$ pair, choose a size-$k$ object from $\mathcal{A}$ and a size-$(n-k)$ object from $\mathcal{B}$. Only finitely many splits $k+(n-k)=n$ are possible.

For sequences, a finite list of $\mathcal{A}$-objects can have length $0,1,2,\dots$, so as a class,
\[
  \SEQ(\mathcal{A}) = \mathcal{E} + \mathcal{A} + \mathcal{A}^2 + \mathcal{A}^3 + \cdots,
\]
where $\mathcal{A}^k$ means a $k$-fold Cartesian product and therefore represents ordered sequences of length $k$. The union is disjoint because the sequence length is part of the data: a 2-term sequence is not the same kind of object as a 3-term sequence.

Why do we exclude size-$0$ objects from $\mathcal{A}$? Because otherwise local finiteness can fail. If $\mathcal{A}$ had even one size-$0$ object, then a size-$1$ sequence could be padded with arbitrarily many zero-size entries, producing infinitely many distinct sequences of total size $1$. By contrast, if every object of $\mathcal{A}$ has size at least $1$, then a sequence of total size $n$ can have length at most $n$. So only finitely many sequence lengths contribute to the coefficient of $z^n$.

Now apply the union and product rules:
\[
  \SEQ(\mathcal{A}) \longleftrightarrow 1 + A(z) + A(z)^2 + A(z)^3 + \cdots.
\]
Because $\mathcal{A}$ has no size-$0$ object, the constant term of $A(z)$ is $0$. So Chapter~1's formal geometric-series identity applies and gives
\[
  1 + A(z) + A(z)^2 + A(z)^3 + \cdots = \frac{1}{1-A(z)}.
\]
Equivalently, the coefficient of $z^n$ on the left counts all sequences of total size $n$, and only finitely many powers $A(z)^k$ can contribute because $k \le n$.
\end{proof}

\begin{remark}
The constructor
\[
  \SEQ_{\ge k}(\mathcal{A})
\]
means sequences of length at least $k$. Structurally, that is
\[
  \SEQ_{\ge k}(\mathcal{A}) = \mathcal{A}^k \times \SEQ(\mathcal{A}),
\]
because one first chooses the required first $k$ entries and then any additional tail. Therefore its OGF is
\[
  \frac{A(z)^k}{1-A(z)}.
\]
In particular,
\[
  \SEQ_{\ge 1}(\mathcal{A}) \longleftrightarrow \frac{A(z)}{1-A(z)}.
\]
\end{remark}

\section{Marking Variables and Bivariate Generating Functions}

A second variable can track a statistic at the same time as size. Let $\chi : \mathcal{C} \to \mathbb{Z}_{\ge 0}$ be a parameter, and define the bivariate OGF
\[
  C(z,u) = \sum_{\gamma \in \mathcal{C}} z^{|\gamma|}u^{\chi(\gamma)} = \sum_{n,k \ge 0} c_{n,k} z^n u^k,
\]
where $c_{n,k}$ counts objects of size $n$ with statistic value $k$. The ordinary OGF is recovered by setting $u=1$.

To speak about an \emph{expected value} at size $n$, we must put a probability distribution on the finite set $\mathcal{C}_n$. In this chapter we use the \emph{uniform distribution} on $\mathcal{C}_n$. Then
\[
  \mathbf{E}[\chi \mid |\gamma|=n]
  = \frac{[z^n] \partial_u C(z,u)\big|_{u=1}}{[z^n] C(z,1)}.
\]
Indeed,
\[
  \partial_u u^{\chi(\gamma)} = \chi(\gamma)u^{\chi(\gamma)-1},
\]
so after setting $u=1$,
\[
  [z^n] \partial_u C(z,u)\big|_{u=1} = \sum_{|\gamma|=n} \chi(\gamma).
\]
Dividing by $|\mathcal{C}_n| = [z^n]C(z,1)$ gives the uniform mean.

A related construction is \emph{pointing}. Suppose each size-$n$ object is built from exactly $n$ atomic positions counted by the variable $z$. Then the pointed class $\mathcal{C}^{\bullet}$ is obtained by choosing one distinguished atom. Every size-$n$ object therefore gives rise to exactly $n$ pointed objects, so
\[
  C^{\bullet}(z) = zC'(z).
\]
For example, a 4-letter word has exactly 4 possible pointings, one for each letter position.

\section{Labeled Classes and EGFs}

For labeled structures, the $n$ atoms carry distinct labels from $\{1,\dots,n\}$. A labeled object is therefore an underlying shape together with a bijection from its $n$ atomic positions to the label set $[n]$. The natural generating function is the \emph{exponential generating function} (EGF)
\[
  \hat{C}(z) = \sum_{n \ge 0} c_n \frac{z^n}{n!}.
\]
The factor $1/n!$ removes the arbitrary naming of the labels and leaves only the structural count.
For example, a labeled object of size $3$ uses the labels $\{1,2,3\}$ exactly once each, assigned to its three atomic positions.

Disjoint union works exactly as before. The crucial new rule is the labeled product. Suppose $\hat{\mathcal{A}}$ and $\hat{\mathcal{B}}$ are labeled classes with EGFs $\hat{A}(z)$ and $\hat{B}(z)$. To build a size-$n$ product object, choose $k$ of the $n$ labels for the $\hat{\mathcal{A}}$ part and give the remaining $n-k$ labels to the $\hat{\mathcal{B}}$ part. Therefore
\[
  n![z^n]\,\hat{A}(z)\hat{B}(z)
  = \sum_{k=0}^n \binom{n}{k} a_k b_{n-k},
\]
which is exactly the number of labeled product objects of size $n$. This is the same binomial-convolution calculation from Chapter~1, now interpreted combinatorially.

If a labeled class has no size-$0$ object, then the sequence rule again has the same formal shape:
\[
  \SEQ(\hat{\mathcal{A}}) \longleftrightarrow \frac{1}{1-\hat{A}(z)}.
\]
The underlying reason is that an ordered list of labeled components still has no symmetry: once the list order is fixed, there is nothing to divide out.

The next labeled constructor is \emph{set}. Here the order of the components is forgotten.

\begin{proposition}[Set constructor in the labeled setting]
If $\hat{\mathcal{A}}$ has no size-$0$ object, then
\[
  \widehat{\SET(\mathcal{A})}(z) = \exp(\hat{A}(z)).
\]
\end{proposition}

\begin{proof}
Let $\SET_k(\mathcal{A})$ denote sets of exactly $k$ labeled $\mathcal{A}$-objects. First pretend the $k$ components are ordered. Then the EGF is $\hat{A}(z)^k$, because we are taking a $k$-fold labeled product. But a set forgets the order of its $k$ components, and there are exactly $k!$ ways to order them. So
\[
  \widehat{\SET_k(\mathcal{A})}(z) = \frac{\hat{A}(z)^k}{k!}.
\]
Now sum over all $k \ge 0$:
\[
  \widehat{\SET(\mathcal{A})}(z)
  = \sum_{k \ge 0} \frac{\hat{A}(z)^k}{k!}
  = \exp(\hat{A}(z)).
\]
The exclusion of size-$0$ objects is again needed to prevent infinitely many zero-size components from appearing.
\end{proof}

As a tiny example, the class of finite sets of labeled atoms is $\SET(\mathcal{Z})$, so its EGF is
\[
  \exp(z).
\]
The coefficient of $z^n/n!$ is $1$, which matches the fact that on a given label set $[n]$ there is exactly one $n$-element set.

Cyclic structures are subtler. The standard labeled rule is
\[
  \widehat{\CYC(\mathcal{A})}(z) = \log\!\frac{1}{1-\hat{A}(z)}.
\]
We record this formula here because it is important and useful, but we will treat it as a black-box rule in this chapter. A full proof belongs to the theory of combinatorial species; see \cite{BergeronLabelleLeRoux1998} or \cite{FlajoletSedgewick2009}.

\section{Worked Examples}

\begin{example}[Binary words]
Let $\mathcal{A}=\{a\}$ and $\mathcal{B}=\{b\}$, each a copy of $\mathcal{Z}$. A binary word is a finite sequence of one-letter choices, so
\[
  \mathcal{W} = \SEQ(\mathcal{A}+\mathcal{B}).
\]
Since each one-letter choice contributes $z$, the class $\mathcal{A}+\mathcal{B}$ has OGF
\[
  z+z = 2z.
\]
Therefore
\[
  W(z) = \frac{1}{1-2z} = \sum_{n \ge 0} 2^n z^n.
\]
So there are $2^n$ binary words of length $n$.
\end{example}

\begin{example}[Plane trees and Catalan numbers]
A plane tree consists of a root together with an ordered list of plane trees attached to that root as children. The children are ordered, so a sequence is exactly the right constructor. If size means number of nodes, then the root contributes one atom and the child list contributes a sequence of smaller trees:

\begin{center}
\begin{tikzpicture}[
  level distance=0.8cm,
  sibling distance=1.2cm,
  every node/.style={circle, draw, fill=black!10, inner sep=1.5pt}
]
  \node (T) {$\mathcal{T}$};
  \node[draw=none, fill=none, right=0.5cm of T] (eq) {$=$};
  \node[right=0.5cm of eq] (root) {};
  \node[draw=none, fill=none, right=0.5cm of root] (times) {$\times$};
  \node[draw=none, fill=none, right=0.5cm of times] (seq) {$\SEQ\Big($};
  \node[right=0.2cm of seq] (T2) {$\mathcal{T}$};
  \node[draw=none, fill=none, right=0.2cm of T2] (close) {$\Big)$};
\end{tikzpicture}
\end{center}

This translates to
\[
  \mathcal{T} = \mathcal{Z} \times \SEQ(\mathcal{T}).
\]
Proposition~\ref{prop:symbolic-ogf} therefore gives
\[
  T(z) = \frac{z}{1-T(z)}.
\]
Rearranging,
\[
  T(z)^2 - T(z) + z = 0.
\]
Solving the quadratic yields
\[
  T(z) = \frac{1 \pm \sqrt{1-4z}}{2}.
\]
Chapter~1 showed that there is a unique formal square root with constant term $1$, and the condition $T(0)=0$ forces the minus sign. So
\[
  T(z) = \frac{1-\sqrt{1-4z}}{2}.
\]
Chapter~1 also computed the coefficients of $\sqrt{1-4z}$, which gives
\[
  [z^n]T(z) = \frac{1}{n}\binom{2n-2}{n-1} \qquad (n \ge 1).
\]
This is the Catalan number $C_{n-1}$. The index shift comes from our size convention: a plane tree with $n$ nodes corresponds to the $(n-1)$st Catalan number because Catalan numbers are often indexed by edges or by internal nodes rather than total nodes.

This example is our first encounter with a square-root expression coming from a symbolic specification. Later chapters will explain when such equations really do force square-root singularities and $n^{-3/2}$ asymptotics; that conclusion is only a preview here.
\end{example}

\begin{example}[Integer compositions]
A positive integer of value $k$ can be modeled as a block of exactly $k$ atoms, because the block size records the integer's value. So the class of positive integers is
\[
  \mathcal{I} = \SEQ_{\ge 1}(\mathcal{Z}),
  \qquad
  I(z) = \frac{z}{1-z}.
\]
A composition is then a nonempty sequence of such blocks:
\[
  \mathcal{K} = \SEQ_{\ge 1}(\mathcal{I}).
\]
Therefore
\[
  K(z) = \frac{I(z)}{1-I(z)} = \frac{z/(1-z)}{1-z/(1-z)} = \frac{z}{1-2z}.
\]
So
\[
  [z^n]K(z) = 2^{n-1} \qquad (n \ge 1).
\]
This matches the direct gaps argument: between $n$ ones there are $n-1$ gaps, and each gap is either cut or not cut.
\end{example}

\begin{example}[Bivariate marking: number of $a$'s in a binary word]
Let $\chi(w)$ count the number of letters $a$ in a binary word $w$. Each letter contributes one factor of $z$ for its size. In addition, the letter $a$ contributes a factor of $u$ while the letter $b$ contributes no extra factor. So a single letter contributes either $zu$ or $z$, and an arbitrary binary word has bivariate OGF
\[
  W(z,u) = \frac{1}{1-(zu+z)} = \frac{1}{1-z(u+1)}.
\]
Now expand the geometric series:
\[
  W(z,u) = \sum_{n \ge 0} z^n (u+1)^n.
\]
Therefore
\[
  [z^n u^k]W(z,u) = [u^k](u+1)^n = \binom{n}{k}.
\]
So among the $2^n$ binary words of length $n$, exactly $\binom{n}{k}$ contain $k$ letters $a$.

Using the uniform distribution on binary words of length $n$, the mean number of $a$'s is
\[
  \frac{[z^n]\partial_u W(z,u)\big|_{u=1}}{[z^n]W(z,1)}
  = \frac{[z^n]\frac{z}{(1-z(u+1))^2}\big|_{u=1}}{2^n}
  = \frac{[z^n]\frac{z}{(1-2z)^2}}{2^n}
  = \frac{n2^{n-1}}{2^n}
  = \frac{n}{2},
\]
as symmetry suggests.
\end{example}

\section{Why the Symbolic Method Matters}

The examples above share a common pattern: write down what an object \emph{is}, check that the decomposition is unique, apply the translation dictionary, and obtain a functional equation for the generating function. That equation may be rational, algebraic, or --- in the labeled setting --- involve familiar transcendental functions such as $\exp$ and $\log$.

At this stage, however, one must be careful not to overstate what has been proved. This chapter establishes the \emph{combinatorial-to-algebraic translation}. It does \emph{not} yet prove the later analytic claims about singularities or asymptotics. For example:
\begin{itemize}
  \item Using only $+$, $\times$, and $\SEQ$ often leads to rational expressions in generating functions, and clearing denominators can produce polynomial equations.
  \item Many context-free grammars later translate into polynomial systems in several generating functions.
  \item Generic square-root singularities often lead to $n^{-3/2}$ factors.
\end{itemize}
But each of these statements requires additional hypotheses and later theorems. In particular, it is \emph{not} true that every algebraic generating function automatically has $n^{-3/2}$ asymptotics: rational functions are algebraic too, and poles lead to different behavior.

Chapter~5 will return to context-free grammars and explain carefully when grammar equations produce algebraic generating functions and when generic square-root asymptotics really appear. The symbolic method is the combinatorial front end of that theory, not the whole theory by itself.

The connection to language models must also be kept precise. In this chapter we are studying \emph{counting} generating functions attached to combinatorial classes. Later chapters will introduce weighted and probabilistic generating functions as separate objects. The symbolic method supplies the counting side of that dictionary, but one should not identify counting OGFs, probability generating functions, and partition functions without saying exactly what has been weighted and how.

\begin{exercise}
Let $\mathcal{P}$ be the class of balanced parenthesis strings (Dyck words). Write a symbolic specification for $\mathcal{P}$ using a decomposition at the first return to height zero, derive the OGF, and verify that
\[
  [z^{2n}]P(z) = C_n.
\]
\end{exercise}

\begin{exercise}
A \emph{unary-binary tree} is either a leaf (one node) or an internal node with exactly one child or exactly two children. Write a symbolic equation for the class $\mathcal{U}$ of unary-binary trees (size = number of nodes), derive the OGF as a root of a quadratic, and identify its radius of convergence.
\end{exercise}

\begin{exercise}[Challenge]
Using the labeled setting and the standard cycle-construction rule recorded above, let $\mathcal{S}$ be the class of permutations and let $\mathcal{D}$ be the class of derangements. Starting from the specification
\[
  \mathcal{D} = \SET(\CYC_{\ge 2}(\mathcal{Z})),
\]
derive the EGF of derangements and confirm that the probability a random permutation of $[n]$ is a derangement tends to $e^{-1}$ as $n \to \infty$.
\end{exercise}

\section*{Further Reading}

\begin{itemize}
  \item \textbf{P. Flajolet and R. Sedgewick}, \emph{Analytic Combinatorics} (Cambridge University Press, 2009) \cite{FlajoletSedgewick2009}, Chapters~I--II --- The definitive treatment of the symbolic method, covering both the unlabeled (OGF) and labeled (EGF) translation dictionaries, with proofs, an extensive worked-example catalogue, and the full derivation of the $\SET$, $\CYC$, and $\MSET$ rules omitted from the present chapter.

  \item \textbf{P. Flajolet and M. Soria}, ``General combinatorial schemas: Gaussian limit distributions and exponential tails,'' \emph{Discrete Mathematics} 114 (1993), 159--180 --- Shows how bivariate generating functions and the symbolic method together yield limit laws (Gaussian distributions of parameters) for broad families of combinatorial structures; a readable illustration of how the bivariate marking machinery of this chapter extends to probability.

  \item \textbf{R.\,P. Stanley}, \emph{Enumerative Combinatorics}, Vol.~2 (Cambridge, 1999) --- Chapter~5 covers trees and the Lagrange inversion formula; Chapter~6 develops algebraic, D-finite, and noncommutative generating functions. Together they extend the symbolic-method viewpoint to more advanced territory.

  \item \textbf{A. Nijenhuis and H.\,S. Wilf}, \emph{Combinatorial Algorithms} (Academic Press, 2nd ed., 1978) \cite{NijenhuisWilf1978} --- Develops algorithms for random generation and enumeration of combinatorial structures (set partitions, permutations, labeled trees); the algorithmic perspective complements the symbolic-method viewpoint of this chapter.

  \item \textbf{F. Bergeron, G. Labelle, and P. Leroux}, \emph{Combinatorial Species and Tree-Like Structures} (Cambridge, 1998) --- The standard reference for the theory of species, the categorical underpinning of the symbolic method; essential for readers who want to understand why the translation dictionary works at the level of bijections rather than just coefficient identities.
\end{itemize}
